% p14-universal-property-of-the-product.tex
%
% Source:   Carl A. Gunter and Dana S. Scott, "Semantic Domains" (1990).
%           ScottDomains/papers/Gunter Scott 1990.pdf
% Physical page: 14          Printed page: 13
% Section:  4.1 "Products".
% Unnamed in the paper - no figure number, no caption.
%
% Introducing sentence (verbatim, printed page 13):
%   "There is another, more pictorial, way of stating these equational
%    properties of the operator \langle\cdot,\cdot\rangle using a commutative
%    diagram. The desired property can be stated in the following manner:
%    given any cpo F and continuous functions f : F \to D and g : F \to E,
%    there is a unique continuous function \langle f,g\rangle which completes
%    the following diagram:"
% and the sentence closing it, immediately below the diagram:
%   "This is referred to as the universal property of the operator \times."
%
% What the development uses it for: it is the universal property of the
% binary product, the equational content of fst \circ \langle f,g\rangle = f,
% snd \circ \langle f,g\rangle = g and \langle fst \circ h, snd \circ h\rangle
% = h, stated as one diagram.
%
% The mediating arrow \langle f,g\rangle is the one the paper draws DASHED -
% it is the arrow asserted to exist and be unique - and that must not be
% flattened to a solid arrow.
%
% Geometry measured from a 600 dpi crop of physical page 14 and mapped to
% centimetres by dividing page pixels by 200: D, D \times E and E sit on one
% vertical 2.4 cm apart, with F 4.2 cm to the right at the height of
% D \times E.
%
% Style contract (r0039) with the orchestrator-settled corrections, plus the
% "obj", "arrow" and "darrow" styles this commutative diagram needs (see
% p08-uniform-fixed-point-operator.tex for the note).

\documentclass[border=6pt]{standalone}
\usepackage{tikz}
\begin{document}
\begin{tikzpicture}[
  x=1cm, y=1cm,
  every node/.style={inner sep=1pt},
  open/.style={circle, draw, fill=white, minimum size=4pt, inner sep=0pt},
  solid/.style={circle, draw, fill=black, minimum size=5pt, inner sep=0pt},
  edge/.style={draw, line width=0.4pt},
  obj/.style={inner sep=3pt},
  arrow/.style={draw, line width=0.4pt, ->},
  darrow/.style={draw, line width=0.4pt, ->, dash pattern=on 4pt off 3pt}]

  \node[obj] (D)  at (0.0, 4.8) {$D$};
  \node[obj] (DE) at (0.0, 2.4) {$D \times E$};
  \node[obj] (E)  at (0.0, 0.0) {$E$};
  \node[obj] (F)  at (4.2, 2.4) {$F$};

  \draw[arrow]  (DE) -- node[left]  {$\mathsf{fst}$} (D);
  \draw[arrow]  (DE) -- node[left]  {$\mathsf{snd}$} (E);
  \draw[arrow]  (F)  -- node[below right] {$f$} (D);
  \draw[arrow]  (F)  -- node[above right] {$g$} (E);
  \draw[darrow] (F)  -- node[above] {$\langle f,g \rangle$} (DE);

\end{tikzpicture}
\end{document}
